\documentclass[11pt]{article}
\usepackage[margin=0.3in,top=0.25in,bottom=0.4in]{geometry}
\makeatletter
\renewcommand{\@maketitle}{\newpage\null\vskip -3.5em\begin{center}{\LARGE \@title \par}\end{center}\par\vskip 0.2em}
\makeatother
\usepackage{amsmath,amssymb}
\usepackage{hyperref}
\usepackage{booktabs}

\title{Why the diffractive/inclusive $D^0$ ratio falls with rapidity in\\
this dipole-model calculation}
\author{}
\date{}

\begin{document}
\maketitle

\section{The puzzle}

Plotting the ratio $d\sigma^{\rm diff}/d\sigma^{\rm incl}$ (this project's
diffractive calculation over \texttt{inclusive-D0-UPC}'s inclusive one) as a
contour in $(y,p_T)$, the ratio \emph{decreases} with increasing rapidity
$y$ in our dipole-model calculation. This is the opposite trend to
Guzey, Innocenti, Sta\'sto \& Strikman, arXiv:2606.05469 (collinear
factorization with nuclear diffractive PDFs, the ``LTA'' shadowing
framework), whose Fig.~7 and Conclusions state explicitly that the
diffractive fraction is \emph{largest} at forward $y$, ``corresponding to
the small values of $x$ in the Pomeron-emitting target nucleus,'' and
decreases at negative rapidity.

Three explanations were checked, two ruled out, one confirmed and
quantified below.

\section{Ruled out: a sign-convention mismatch}

Both projects define $x\propto e^{-y}$ (smaller $x$ at larger $y$): this
project's photon plus-momentum threshold $pp=(m_T/\sqrt2)\,e^{y}$ grows
with $y$, and the sibling inclusive project's own
$x_{bj}=(m_{T0}/\sqrt{s})\,e^{-y}$ shrinks with $y$ -- the same direction.
arXiv:2606.05469's own stated convention (Sec.~7: forward $y\Rightarrow$
small $x$ in the target) matches this too. So all three sources agree on
what ``positive $y$'' means; this is not a mislabeled axis.

\section{Ruled out: the diffractive integrand's own kinematic ceiling}

The diffractive integrand enforces $qp\ge pp$, capping the integration at
$qp_{\max}=800$~GeV (\texttt{jac\_qp = qpmax - pp; if (jac\_qp<=0) return
0;}). At $p_T=1$~GeV, $pp$ only reaches $\sim 26$~GeV by $y=3$ -- nowhere
near the 800~GeV ceiling, so this constraint never actually activates in
the plotted range and cannot explain the suppression.

Also ruled out: \texttt{inclusive-D0-UPC} lacking an analogous
kinematic threshold. It does have one -- \texttt{inclusive\_spectrum.cpp}
recomputes the identical $pp$ formula locally per sample
(\texttt{pp\_zh}), with the same \texttt{qpmax} ceiling. The unused
\texttt{param.qpmin} field in its \texttt{main.cpp} is dead code, but the
physical constraint itself is present and applied the same way in both
calculations.

\section{Confirmed: the dipole amplitude does grow -- but the
\emph{flux-weighted} $x$ diffraction effectively probes doesn't shrink
as fast as $x_{bj}$}

\subsection{$N(r,x)$ does grow toward small $x$}

Querying this project's own Pb dipole amplitude directly
(\texttt{AmplitudeLib::N(r,x)}, central impact parameter,
$r=1/m_c\approx0.667$~GeV$^{-1}$, the charm-mass dipole size):

\begin{center}
\begin{tabular}{lr}
\toprule
$x$ & $N(r{=}1/m_c,\,x)$ \\
\midrule
$10^{-2}$ (or larger, clamped) & 0.093 \\
$10^{-3}$ & 0.153 \\
$10^{-4}$ & 0.228 \\
$10^{-5}$ & 0.321 \\
\bottomrule
\end{tabular}
\end{center}

$N$ more than triples over this range and stays well below 1 (not
saturated). For small $N$, $d(N^2)/N^2 = 2\,dN/N$ -- i.e.\ $N^2$
(diffractive-like scaling) grows \emph{faster in relative terms} than
$N$ itself as $x$ falls. Naively this predicts the diffractive/inclusive
ratio should \emph{increase} with $y$, matching arXiv:2606.05469's
finding and contradicting what we compute. The dipole amplitude itself is
behaving exactly as expected; the discrepancy must be elsewhere.

\subsection{The diffractive integral doesn't sample $x_{\mathbb P}$ at
$x_{bj}$}

The diffractive kernel evaluates $N(r,x_{\mathbb P})$ with
\[
x_{\mathbb P} = \frac{M_{q\bar q}^2}{w^2}, \qquad
w^2 = \sqrt2\,qp\,\sqrt{s}, \qquad
M_{q\bar q}^2 = \frac{p_T^2+m_c^2}{z_1 z_2}, \qquad
z_1=\frac{pp}{qp},\ z_2=1-z_1,
\]
integrated over $qp\in[pp,\,qp_{\max}]$ and weighted by the photon flux.
At the upper end of that range ($qp\to\infty$, $z_1\to0$),
$x_{\mathbb P}\to m_T e^{-y}/\sqrt{s} = x_{bj}$ exactly -- the two
variables coincide there. But at the \emph{lower} end ($qp\to pp$,
$z_1\to1$), $M_{q\bar q}^2\to\infty$ and $x_{\mathbb P}$
\emph{diverges}. The photon flux is steeply peaked at low $qp$ -- right
next to that divergent lower edge -- so the flux-weighted average
$x_{\mathbb P}$ the integral actually samples is pulled well above
$x_{bj}$, not equal to it.

Computed directly (\texttt{photon\_flux()}, central impact parameter,
$p_T=1$~GeV, $\sqrt{s_{NN}}=5.36$~TeV):

\begin{center}
\begin{tabular}{lrrr}
\toprule
$y$ & $x_{bj}$ (inclusive) & $\langle x_{\mathbb P}\rangle_{\rm flux}$ (diffractive) & ratio \\
\midrule
$-2.0$ & $2.49\times10^{-3}$ & $6.96\times10^{-3}$ & 2.8 \\
$0.0$  & $3.36\times10^{-4}$ & $1.32\times10^{-3}$ & 3.9 \\
$1.5$  & $7.50\times10^{-5}$ & $4.55\times10^{-4}$ & 6.1 \\
$3.0$  & $1.67\times10^{-5}$ & $2.15\times10^{-4}$ & 12.8 \\
\bottomrule
\end{tabular}
\end{center}

The offset is not constant: it \emph{widens} from $\sim2.8\times$ at
$y=-2$ to $\sim12.8\times$ at $y=3$ -- roughly a factor of 4.6 in
relative terms across the plotted range. As $y$ increases, the
diffractive calculation effectively probes a progressively \emph{larger}
$x$ than the $x_{bj}$ that governs inclusive production at the same
nominal $y$, starving diffraction of exactly the low-$x$ enhancement
that's driving inclusive's growth. This directly explains the observed
sign of the effect and is large enough to plausibly account for its
size, though that has not been checked point-by-point against the
full numerical ratio.

\section{Conclusion}

The falling diffractive/inclusive ratio with $y$ is not a sign-convention
bug or an artificial kinematic cutoff; it traces to how the photon flux
weights the $qp$ integration relative to the $y$-dependent lower bound
$pp$, which systematically shifts the effective $x_{\mathbb P}$ sampled
by the diffractive calculation above the $x_{bj}$ relevant for inclusive,
increasingly so at large $y$. This is a real feature of \emph{this
formula} for how $x_{\mathbb P}$ enters the diffractive dipole
amplitude -- it is not obviously a coding error, but it is a legitimate
open question whether this is the physically correct way to define/average
$x_{\mathbb P}$ in this kind of calculation, or whether the disagreement
with arXiv:2606.05469 simply reflects the two frameworks (dipole/CGC vs.\
collinear factorization + nuclear diffractive PDFs) genuinely predicting
different rapidity dependence. Resolving that fully would need literature
comparison with other dipole-model diffractive-DIS treatments of
$x_{\mathbb P}$, beyond what code inspection alone can settle.

\end{document}
